\clearpage
\section{Nematic Order}
\begin{colbox}
    \begin{enumerate}
        \item The nematic tensor in the fully ordered state is given by \(Q_{\alpha\beta} = n_\alpha n_\beta-\delta_{\alpha\beta}/d\) where \(n_\alpha\) is the director field and \(d\) is the dimension. Show that in 2D \begin{equation}
            \vec{Q} = \frac{1}{2}\begin{pmatrix}
                \cos 2\theta & \sin 2\theta \\
                \sin 2\theta & -\cos 2\theta
            \end{pmatrix}
        \end{equation}
        where \(\vec{n} = \ins{\cos\theta,\sin\theta}\).
        \item Consider the Frank free energy
        \begin{equation}
            F = \int\dd \vec{r}\left[\frac{1}{2}K_1\ins{\nabla\cdot\vec{n}}^2+\frac{1}{2}K_2{\left[\vec{n}\cdot\ins{\nabla\times\vec{n}}\right]}^2+\frac{1}{2}K_3{\left[\vec{n}\times\ins{\nabla\times\vec{n}}\right]}^2\right]
        \end{equation} 
        Show that within the one-constant approximation \((K_i=K)\), the free energy reduces to
        \begin{equation}
            F = \int \dd \vec{r} \frac{1}{2}K\del_\alpha n_\beta\del_\alpha n_\beta + F_b
        \end{equation}
        where \(F_b\) is a boundary term, called saddle-splay. In 2D, show that this free energy can also be written as \(F = \int\dd \vec{r}\frac{1}{2}K\ins{\nabla\theta}^2\).
    \end{enumerate}
\end{colbox}
\subsection{Nematic Tensor}
By a direction calculation, in 2D
\begin{align}
    Q_{\alpha\beta} &= n_\alpha n_\beta-\delta_{\alpha\beta}/2 \\
    &= \begin{pmatrix}
        \cos^2\theta-\frac{1}{2} & \cos\theta\sin\theta \\
        \sin\theta\cos\theta & \sin^2\theta-\frac{1}{2}
    \end{pmatrix}\\
    &= \frac{1}{2}\begin{pmatrix}
        2\cos^2\theta-1 & 2\cos\theta\sin\theta \\
        2\sin\theta\cos\theta & 2\sin^2\theta - 1
    \end{pmatrix} \\
    &= \frac{1}{2}\begin{pmatrix}
        2\cos^2\theta-\sin^2\theta-\cos^2\theta & 2\cos\theta\sin\theta \\
        2\sin\theta\cos\theta & 2\sin^2\theta - \sin^2\theta-\cos^2\theta
    \end{pmatrix}\\
    &= \frac{1}{2}\begin{pmatrix}
        \cos^2\theta-\sin^2\theta & 2\cos\theta\sin\theta \\
        2\sin\theta\cos\theta & \sin^2\theta -\cos^2\theta
    \end{pmatrix} \\
    &= \frac{1}{2}\begin{pmatrix}
        \cos 2\theta & \sin 2\theta \\
        \sin 2 \theta & -\cos 2\theta
    \end{pmatrix}
\end{align}
where the last transition used the trig identities
\begin{align}
    \sin 2\theta = 2\sin\theta\cos\theta && \cos 2\theta = \cos^2\theta - \sin^2\theta
\end{align}
\subsection{Frank Free Energy}
Since we are working within the one-constant approximation, we can immediately simplify the free energy to
\begin{equation}
    F = \int\dd \vec{r}\frac{K}{2}\left[\ins{\nabla\cdot\vec{n}}^2+{\left[\vec{n}\cdot\ins{\nabla\times\vec{n}}\right]}^2+{\left[\vec{n}\times\ins{\nabla\times\vec{n}}\right]}^2\right]
\end{equation}
if we define for simplicity \(\vec{a} = \nabla\times\vec{n}\) then
\begin{equation}
    F = \int\dd \vec{r}\frac{K}{2}\left[\ins{\nabla\cdot\vec{n}}^2+{\left[\vec{n}\cdot\vec{a}\right]}^2+{\left[\vec{n}\times\vec{a}\right]}^2\right]
\end{equation}
but
\begin{align}
    \ins{\vec{n}\times\vec{a}}^2 &= \ins{\epsilon_{ijk}n_j a_k}\ins{\epsilon_{lmn}n_m a_n} = \epsilon_{ijk}\epsilon_{lmn}n_j a_k n_m a_n \\
    &= \ins{\delta_{jm}\delta_{kn}-\delta_{jn}\delta_{mk}}n_j n_m a_k a_n \\
    &= n_j n_j a_k a_k - n_j n_k a_j a_k \\
    &= n^2 a^2 - \ins{\vec{n}\cdot\vec{a}}^2
\end{align}
therefore
\begin{equation}
    F = \int\dd \vec{r}\frac{K}{2}\left[\ins{\nabla\cdot\vec{n}}^2 + n^2 a^2\right]
\end{equation}
but since \vec{n} is a unit vector
\begin{equation}
    F = \int\dd \vec{r}\frac{K}{2}\left[\ins{\nabla\cdot\vec{n}}^2 + \ins{\nabla\times\vec{n}}^2 \right]
\end{equation}
converting to index notation
\begin{align}
    F &= \int\dd \vec{r}\frac{K}{2}\left[\del_i n_i \del_j n_j + \del_j n_k \del_j n_k - \del_j n_k \del_k n_j \right] \\
    &= \int\dd \vec{r}\frac{K}{2}\left[\del_i n_j \del_i n_j - \del_i\ins{n_j \del_j n_i - n_i \del_j n_j}  \right]
\end{align}
if we identify \(F_b = - \del_i\ins{n_j \del_j n_i - n_i \del_j n_j}\) we found what's needed. If we now assume we are in 2D then
\begin{equation}
    \vec{n} = \ins{\cos\theta, \sin\theta}
\end{equation}
therefore
\begin{equation}
    \del_\alpha \vec{n} = \del_\theta \vec{n} \del_\alpha \theta = \ins{-\sin\theta,\cos\theta} \del_\alpha \theta
\end{equation}
and thus
\begin{equation}
    \ins{\del_\alpha \vec{n}}^2 = \ins{\sin^2\theta + \cos^2\theta}\ins{\del_\alpha \theta}^2 = \ins{\del_\alpha \theta}^2 = \ins{\nabla\theta}^2
\end{equation}
which gives us
\begin{equation}
    \int\dd \vec{r}\frac{K}{2}\ins{\nabla\theta}^2
\end{equation}